\documentclass[UTF8,11pt]{ctexart}

\usepackage[a4paper,margin=2.3cm]{geometry}
\usepackage{amsmath,amssymb,bm}
\usepackage{booktabs,array}
\usepackage{xcolor}
\usepackage{enumitem}
\usepackage{listings}
\lstset{
  basicstyle=\ttfamily\footnotesize,
  keywordstyle=\color{blue!70!black},
  commentstyle=\color{gray!60!black}\itshape,
  stringstyle=\color{teal!70!black},
  breaklines=true, showstringspaces=false,
  columns=fullflexible, keepspaces=true,
  frame=single, rulecolor=\color{gray!40}, framesep=4pt,
  xleftmargin=4pt, xrightmargin=2pt, aboveskip=4pt, belowskip=2pt,
  language=Python, escapeinside={(*}{*)},
}
% one-line tag above a code block: filename / function it comes from
\newcommand{\codefrom}[1]{\vspace{1pt}\noindent{\small\color{gray!55!black}\textsf{代码 · #1}}\par\nopagebreak}
\usepackage{hyperref}
\hypersetup{colorlinks=true,linkcolor=blue,urlcolor=blue,citecolor=blue}

% --- shared macros (consistent with docs/40) ---
\newcommand{\Vp}{V_\mathrm{p}}
\newcommand{\Sw}{S_\mathrm{w}}
\newcommand{\Sp}{S_\mathrm{p}}
\newcommand{\Aflux}{A_\mathrm{flux}}
\newcommand{\Lax}{L_\mathrm{ax}}
\newcommand{\nfp}{N_\mathrm{fp}}
\newcommand{\Pfus}{P_\mathrm{fus}}
\newcommand{\Pbrem}{P_\mathrm{brem}}
\newcommand{\Pcycl}{P_\mathrm{cycl}}
\newcommand{\Pheat}{P_\mathrm{heat}}
\newcommand{\Pwall}{P_\mathrm{wall}}
\newcommand{\Eth}{E_\mathrm{th}}
\newcommand{\iotabar}{\bar{\iota}}
\newcommand{\nphi}{n_\varphi}
\newcommand{\etabar}{\bar\eta}
\newcommand{\Nhat}{\hat{\bm n}}
\newcommand{\Bhat}{\hat{\bm b}}
\newcommand{\That}{\hat{\bm t}}

\title{\textbf{PolyFusion 仿星器位形：完全报告}\\[2pt]
\large 物理模型 · 几何 · 可视化 · 预设库 · POPCON 与性能}
\author{PolyFusion release —— 仿星器位形系列报告（配 docs/39、docs/40）}
\date{\today}

\begin{document}
\maketitle

\begin{abstract}
本报告把 PolyFusion 0-D 程序里\textbf{整个仿星器位形}从头讲透：它\textbf{算什么物理}（0-D 功率平衡）、\textbf{怎么描述几何}（近轴展开 + 真机傅里叶边界）、\textbf{几何如何进功率账}（哪个量、什么幂次）、\textbf{怎么画截面}（统一嵌套磁面 + 相位滑块）、\textbf{内置哪些位形}（9 个预设），以及 \textbf{POPCON 扫描怎么算、怎么被加速}。目标是让懂等离子体物理、但没读过这套代码的人，能完整地"听懂"每一步在做什么、为什么这么做、近似在哪里。核心观念贯穿全文：功率公式只含\emph{标量} $(a_\mathrm{vol},R_0,\Vp,\Sw,\Lax,\iotabar)$，无显式"形状"项；\textbf{但} $\Vp,\Sw$ 是\emph{前端所画 $\rho{=}1$ 边界的精确积分或实测覆写}，所以对概念堆\textbf{边界形状经 $\Vp/\Sw$ 真的进功率}（改塑形 $\etabar$ 可让 $\Pfus$ 摆动 $\sim5\%$），真机则用实测 $\Vp/\Sw/\iotabar$ 覆写、形状主要供显示和几何诊断。前后端共享同一条边界；一阶截面积诊断为 $\pi a^2$，求解器输出有效 $\Aflux=\Vp/\Lax$。纯装饰、不进功率的只有内面嵌套画法与相位帧。
\end{abstract}

\tableofcontents
\vspace{6pt}

%======================================================================
\section{引言：这套代码在做什么}

\paragraph{仿星器 vs 托卡马克。}
托卡马克靠等离子体里流过的大电流产生极向磁场，把磁力线扭成嵌套的环面磁面；仿星器\textbf{不靠等离子体电流}，而是用\textbf{外部三维线圈}把磁面直接"拧"出来。代价是位形天生三维、截面随环向角 $\varphi$ 变化（W7-X 的 bean$\to$triangle，LHD 的旋转椭圆）；好处是没有电流驱动、原则上可稳态、不会有电流驱动不稳定性。

\paragraph{0-D 系统码。}
PolyFusion 是\textbf{零维（0-D）功率平衡}程序：它\emph{不}解 MHD 平衡或输运方程，而是把一台装置压缩成一组标量（体积、壁面积、磁轴长度、旋转变换、温度密度剖面指数……），用代数公式估算\textbf{聚变功率、辐射、约束时间、增益 $Q$、$\beta$、壁负荷}等，用来快速扫参数、画运行空间图（POPCON）。它和同程序里的托卡马克、磁镜、FRC、偶极场共用同一套功率平衡内核，区别只在\textbf{几何标量怎么取}和\textbf{约束定标用哪条经验律}。

\paragraph{本报告范围。}
只讲\textbf{仿星器位形}：第 \ref{sec:power} 节物理内核，第 \ref{sec:nae}--\ref{sec:integral} 节几何（近轴、真机边界、体积积分），第 \ref{sec:account} 节几何进功率，第 \ref{sec:view} 节截面可视化，第 \ref{sec:presets} 节 9 个预设，第 \ref{sec:popcon} 节 POPCON 与性能，第 \ref{sec:consistency} 节一致性与局限。代码在 \texttt{polyfusion/configs/stellarator.py}、\texttt{polyfusion/nearaxis.py}、\texttt{polyfusion/presets/stellarator.json}。\textbf{本报告在关键公式旁贴上对应的真实代码片段}（摘自当前 HEAD），便于公式$\leftrightarrow$实现逐行对照；片段保留原英文注释。

%======================================================================
\section{物理内核：0-D 功率平衡}\label{sec:power}

一切从\textbf{剖面假设}开始。径向坐标 $x=r/a\in[0,1]$，密度与温度取
\begin{equation}
n(x)=n_0(1-x^2)^{S_n},\qquad T(x)=T_0(1-x^2)^{S_T},
\end{equation}
$S_n,S_T$ 是剖面峰化指数（输入）。所有体积量都是这两条剖面在 $\Vp$ 上的积分。

\subsection{组分与有效电荷}
给定离子混合（主离子份额 $f_1$、氦灰 $f_\mathrm{He}$、杂质 $f_\mathrm{imp}$、杂质电荷 $Z_\mathrm{imp}$），由准中性算出中心电子密度 $n_{e0}$、有效电荷
\[
Z_\mathrm{eff}=\frac{\sum_s n_{s0}Z_s^2}{n_{e0}},
\]
和平均质量数 $M$。这步与托卡马克逐字相同（已 L3 验证）。

\subsection{聚变功率：反应性体积积分}
Maxwellian 反应率 $\langle\sigma v\rangle(T)$ 由 \texttt{reactivity} 给出（D-T、D-D、D-He3、p-B11 等，由 \texttt{icase} 选）。把温度剖面代入、对 $x$ 积分得到\textbf{剖面平均反应性因子}
\begin{equation}
\Phi=f_\sigma\cdot 2\!\int_0^1 (1-x^2)^{2S_n}\,\langle\sigma v\rangle\!\big(T(x)\big)\,x\,dx,
\end{equation}
于是
\begin{equation}
\boxed{\;\Pfus=\frac{Y}{1+\delta_{12}}\,n_{10}n_{20}\,\Phi\,\Vp\;}\qquad(\text{乘 }10^{-6}\text{ 转 MW}),
\end{equation}
$Y$ 是单次反应释能、$\delta_{12}$ 处理同种反应物的双计数。中子功率 $P_n=\Pfus(1-f_\mathrm{ion})$，$f_\mathrm{ion}$ 是带电产物（如 $\alpha$）的能量份额。\textbf{关键}：$\Pfus\propto\Vp$——体积是最强的几何杠杆。

\codefrom{\texttt{stellarator.py} · \texttt{solve\_stellarator}}
\begin{lstlisting}
x = np.linspace(0.0, 1.0, 101); dx = x[1] - x[0]
Tx = Ti0 * (1 - x**2) ** ST                      # temperature profile T(x)
sgv = reactivity(Tx, icase)                      # <sigma v>(T(x))
Phi = fsig * 2 * float(np.sum((1 - x**2)**(2*Sn) * sgv * x * dx))   # profile-avg
Pfus = rx["Y"] / (1 + d12) * n10 * n20 * Phi * Vp * 1e-6           # proportional Vp
Pn = Pfus * (1 - rx["fion"])
\end{lstlisting}

\subsection{辐射损失}
\begin{itemize}[leftmargin=1.4em]
\item \textbf{轫致辐射} $\Pbrem\propto n_e^2\sqrt{T_e}\,\Vp$（含相对论修正、$Z_\mathrm{eff}$ 加权），$\propto\Vp$。
\item \textbf{回旋辐射} $\Pcycl\propto B_0^{2.5}\,T_\mathrm{eff}^{2.5}\,n_\mathrm{eff}^{0.5}(1-R_w)^{0.5}\,a_\mathrm{vol}^{-0.5}\,\Vp$，$R_w$ 是壁反射率。注意它\textbf{显含体积等效小半径} $a_\mathrm{vol}^{-1/2}$——回旋辐射对装置尺寸敏感。
\item \textbf{杂质线辐射} $P_\mathrm{line}$（Mavrin 拟合，可选）。
\end{itemize}

\codefrom{\texttt{stellarator.py} · \texttt{solve\_stellarator}（辐射）}
\begin{lstlisting}
Pbrem = (5.34e-37 * ne0**2 * math.sqrt(Te0)
         * (Zeff * (1 / (1 + 2*Sn + 0.5*ST))
            + 0.7936 / (1 + 2*Sn + 1.5*ST) * (Te0 / MEC2)
            + 1.874 / (1 + 2*Sn + 2.5*ST) * (Te0 / MEC2)**2
            + 3 / math.sqrt(2) / (1 + 2*Sn + 1.5*ST) * (Te0 / MEC2))
         * 1e-6 * Vp)
neff = ne0 / 1e20 / (1 + Sn)
Teff = Te0 * float(np.sum((1 - x**2) ** ST)) * dx
Pcycl = (4.14e-7 * neff**0.5 * Teff**2.5 * B0**2.5 * (1 - Rw)**0.5
         * a_vol**-0.5 * (1 + 2.5 * Teff / 511) * Vp)   # volume-equivalent radius
\end{lstlisting}

\subsection{储能、加热、增益}
储能与输运损失
\begin{equation}
\Eth=\tfrac32\,(n_{i0}T_{i0}+n_{e0}T_{e0})\,\frac{\Vp}{1+S_n+S_T},\qquad
P_\mathrm{tr}=\Eth/\tau_E .
\end{equation}
净外加热由功率守恒定出
\begin{equation}
\Pheat=\Pcycl+\Pbrem+P_\mathrm{line}+\Eth/\tau_E-f_\mathrm{ion}\Pfus,
\end{equation}
增益 $Q=\Pfus/\Pheat$（$\Pheat\le0$ 即\textbf{点火}）。环向比压
\[
\beta_t=\frac{2\mu_0(n_{i0}T_{i0}+n_{e0}T_{e0})}{B_0^2(1+S_n+S_T)} .
\]

\codefrom{\texttt{stellarator.py} · \texttt{solve\_stellarator}（储能/加热/$\beta$/壁负荷）}
\begin{lstlisting}
Eth = 1.5 * (ni0*Ti0 + ne0*Te0) * 1e3 * QE / (1 + Sn + ST) * Vp * 1e-6
Pth = Eth / tauE
Pheat = Pcycl + Pbrem + P_line + Pth - rx["fion"] * Pfus
Qfus = Pfus / Pheat if Pheat > 0 else 1000.0      # Pheat<=0 -> ignited
betaT = 2 * MU0 * (ni0*Ti0 + ne0*Te0) * 1e3 * QE / B0**2 / (1 + Sn + ST)
Pwall = (Pfus + Pheat) / Sw                        # the only Sw user
\end{lstlisting}

\subsection{约束闭合：ISS04 + Sudo}
仿星器\textbf{不用}托卡马克的 H98，而用国际仿星器定标 \textbf{ISS04}：
\begin{equation}
\boxed{\;\tau_\mathrm{ISS04}=0.134\,f_\mathrm{ren}\,a_\mathrm{vol}^{2.28}R_0^{0.64}P_L^{-0.61}
\Big(\tfrac{\bar n}{10^{19}}\Big)^{0.54}B_0^{0.84}\,\iotabar^{0.41}\;}
\end{equation}
$P_L=f_\mathrm{ion}\Pfus+\Pheat$ 是损失功率，$f_\mathrm{ren}$ 是装置归一化因子，$H_\mathrm{ISS04}=\tau_E/\tau_\mathrm{ISS04}$ 是约束品质（实验在 $1$--$1.4$）。\textbf{旋转变换 $\iotabar$ 仅在此处以 $0.41$ 次幂进入}——这是几何（或实测）$\iotabar$ 对功率账的唯一通道。密度上限用 Sudo：
\[
n_\mathrm{Sudo}=0.25\sqrt{P_L B_0/(a_\mathrm{vol}^2 R_0)}\times10^{20}.
\]
两个自洽回路：$\tau_E=0$ 时从 ISS04 预测求解 $\tau_E$（令 $H_\mathrm{ISS04}=H_\mathrm{fac}$）；$f_T=0$ 时从电子通道平衡自洽解 $T_e/T_i$。

\codefrom{\texttt{stellarator.py} · \texttt{solve\_stellarator}（ISS04 + Sudo）}
\begin{lstlisting}
PL = rx["fion"] * Pfus + Pheat                     # loss power
tau_ISS04 = (0.134 * f_ren * a_vol**2.28 * R0**0.64 * PL**-0.61
             * (nbar / 1e19)**0.54 * B0**0.84 * iota_used**0.41)  # iota^0.41
H_ISS04 = tauE / tau_ISS04
n_Sudo = 0.25 * math.sqrt(PL * B0 / (a_vol**2 * R0)) * 1e20
\end{lstlisting}

\paragraph{小结。}功率账只认一组\textbf{标量}：$\Vp$（最强，进 $\Pfus,\Pbrem,\Pcycl,\Eth$）、$\Sw$（只进壁负荷 $\Pwall=(\Pfus+\Pheat)/\Sw$）、$a_\mathrm{vol},R_0,\iotabar$（进 Pcycl/ISS04/Sudo）。功率公式里没有"形状"这一项——\textbf{但这不等于形状不进功率}：第 \ref{sec:integral} 节会看到 $\Vp,\Sw$ 是\textbf{所画 $\rho{=}1$ 边界的精确积分或真机实测覆写}，所以\textbf{边界形状经 $\Vp/\Sw$ 进功率}（概念堆；真机用实测覆写）。下面几节就讲这些标量从哪来、形状如何经它们影响功率。

%======================================================================
\section{几何模型 I：近轴展开（概念堆）}\label{sec:nae}

\paragraph{动机。}
怎么用\emph{少数几个数}描述一个三维仿星器位形？现代优化仿星器（QA/QH/QI）几乎都用\textbf{近轴展开}（near-axis / Garren--Boozer，Landreman--Sengupta--Plunk 形式）：把磁面沿磁轴做小半径 $r$ 的泰勒展开。一阶给\textbf{旋转椭圆}，二阶弯成\textbf{bean/crescent}。它只需要：磁轴的傅里叶系数 $\{R_{c,n}\},\{Z_{s,n}\}$、一个塑形标量 $\etabar$、场周期数 $\nfp$，就能解析地、自洽地给出旋转变换、拉长比、截面形状。PolyFusion 的\textbf{概念堆}预设就用这条路。

\subsection{磁轴与 Frenet 标架}
磁轴是三维傅里叶闭曲线
\begin{equation}
R_0(\varphi)=\sum_n R_{c,n}\cos(n\nfp\varphi),\qquad
Z_0(\varphi)=\sum_n Z_{s,n}\sin(n\nfp\varphi).
\end{equation}
默认螺旋轴 $R_{c}=[R_0,\delta_h],\,Z_{s}=[0,-\delta_h]$（$\delta_h$ 为螺旋偏移）；概念堆也可给显式多谐波 $rc/zs$。沿轴建 Frenet 标架（法向 $\Nhat$、副法向 $\Bhat$、切向 $\That$），算出\textbf{曲率} $\kappa(\varphi)$ 和\textbf{挠率} $\tau(\varphi)$。代码用解析傅里叶导数 + 谱微分矩阵，所有量在一个场周期 $[0,2\pi/\nfp)$ 的 $\nphi$ 点网格上。

\paragraph{三种用户入口，但只有两条几何主路径。}
从 UI 和预设看，仿星器几何有三种入口：第一，初学者/概念扫描使用 $R_0,a,\nfp,\delta_h,\etabar$，由代码自动生成默认螺旋磁轴；第二，pyQSC / near-axis 用户使用显式 $rc/zs+\etabar$，直接给磁轴 Fourier；第三，DESC/VMEC 边界用户和真机预设使用 \texttt{shape.R/Z} 的边界 Fourier 三元组。前两者\textbf{同属近轴路径}，都会进入 \texttt{solve\_near\_axis()}；第三者走\textbf{真机边界路径}，由 \texttt{\_machine\_boundary\_outlines()} / \texttt{\_shape\_boundary\_fn()} 直接求边界。

\begin{center}
\renewcommand{\arraystretch}{1.25}
\begin{tabular}{@{}llll@{}}
\toprule
入口 & 适合用户 & 数据形式 & 几何含义 \\
\midrule
$\delta_h+\etabar$ & 初学者、概念堆扫参 & 标量 & 默认螺旋磁轴 + 近轴 r2 磁面 \\
$rc/zs+\etabar$ & pyQSC / Garren--Boozer 近轴用户 & 一维数组 & 显式\textbf{磁轴} Fourier + 近轴 r2 磁面 \\
\texttt{shape.R/Z} & DESC/VMEC 边界用户、真机用户 & $[m,n,c]$ 三元组列表 & 二维\textbf{边界面} Fourier \\
\bottomrule
\end{tabular}
\end{center}

因此 precise-QA / precise-QH 的 $rc/zs$ 与 W7-X 的 \texttt{shape.R/Z} 不是同一种 Fourier：前者是一条磁轴曲线 $R_0(\varphi),Z_0(\varphi)$；后者是完整边界面 $R(\theta,\varphi),Z(\theta,\varphi)$ 的截断谐波。UI 用三档几何输入模式开关把二者语义分开：\textbf{默认近轴}只露 $\delta_h,\etabar$；\textbf{磁轴 Fourier}只露 $rc/zs+\etabar$；\textbf{边界 Fourier}只露边界 $R/Z$ Fourier 与真机覆写 $\iotabar,\Vp^*,\Sw^*$。内部仍以隐藏 \texttt{shape} 对象透传，避免编辑普通数值后丢失真机边界。

\subsection{$\sigma$ 方程与旋转变换}
一阶磁面的形状由 Garren--Boozer 的\textbf{周期 $\sigma$ 方程}决定，它是关于 $\sigma(\varphi)$ 和轴上旋转变换 $\iota$ 的非线性常微分方程，含\textbf{挠率贡献}。代码用 Newton 迭代解出 $(\iota,\sigma)$（对齐 pyQSC 到机器精度 $<2\times10^{-12}$）。直观理解：$\iotabar$ = 磁力线绕环一圈、极向转过的圈数；它由\textbf{轴的扭转（挠率）}和\textbf{塑形}共同决定——所以螺旋偏移 $\delta_h$ 越大、$\iotabar$ 越不同（旋转椭圆模型对 $\delta_h$ 是盲的，这正是本程序弃用它的原因）。

\codefrom{\texttt{nearaxis.py} · \texttt{solve\_near\_axis}（$\sigma$ 方程 Newton 迭代）}
\begin{lstlisting}
x = np.zeros(nphi)                       # x[0] = iota, x[1:] = sigma(phi)
for _ in range(newton_maxit):
    r = residual(x)
    if np.max(np.abs(r)) < newton_tol:   # newton_tol = 1e-12
        break
    x = x - np.linalg.solve(jacobian(x), r)   # <- the BLAS-heavy solve (sec 9)
iota = float(x[0]); sigma = x.copy(); sigma[0] = 0.0
\end{lstlisting}

\subsection{一阶截面：旋转椭圆与拉长比}
解出 $\sigma$ 后，一阶 Frenet 系数
\begin{equation}
X_{1c}=\frac{\etabar}{\kappa},\quad Y_{1s}=\frac{\kappa}{\etabar},\quad Y_{1c}=\frac{\kappa\,\sigma}{\etabar},
\end{equation}
截面是椭圆，拉长比
\begin{equation}
e=\frac{p+\sqrt{p^2-4}}{2},\quad p=X_{1c}^2+Y_{1s}^2+Y_{1c}^2 ,
\end{equation}
随 $\etabar$ 减小而增大（$\sim(\kappa/\etabar)^2$）。\textbf{重要}：一阶面积 $\propto X_{1c}Y_{1s}=1$（面积守恒）——拉长比变、面积不变，这就是一阶截面积诊断 $\pi a^2$ 的根据；求解器另输出有效 $\Aflux=\Vp/\Lax$。

\codefrom{\texttt{nearaxis.py} · \texttt{solve\_near\_axis}（一阶系数与拉长比）}
\begin{lstlisting}
X1c = etabar / curvature
Y1s = curvature / etabar
Y1c = curvature * sigma / etabar
p = X1c**2 + Y1s**2 + Y1c**2
q = -X1c * Y1s                              # = -1 (area-preserving)
elongation = (p + np.sqrt(p*p - 4*q*q)) / (2 * np.abs(q))
\end{lstlisting}

\subsection{二阶截面：bean/crescent}
二阶（$r^2$）系数（LSP Part 2，真空、仿星器对称特例，移植 pyQSC \texttt{calculate\_r2}，验证 $<10^{-9}$）给出
\[
P(\theta)=\text{axis}+X\Nhat+Y\Bhat+Z\That,\quad
X=rX_{1c}\cos\theta+r^2(X_{20}+X_{2c}\cos2\theta+X_{2s}\sin2\theta),\ \dots
\]
$r^2$ 项把一阶椭圆\textbf{弯成 bean/crescent}——真机的特征截面。$B_{2c}$ 是二阶场强塑形输入（默认 $0$）。\textbf{注意}：$r^2$ 展开只在小 $r$ 有效；到反应堆小半径 $a$，$r^2$ 项会超过 $r^1$ 项使面自交——这决定了显示要限幅（第 \ref{sec:view} 节）。

%======================================================================
\section{几何模型 II：真机的傅里叶边界谐波}\label{sec:machine}

\paragraph{为什么近轴不够。}
单谐波近轴\textbf{无法}表示真实实验装置：W7-X 是准等动（bean$\to$triangle）、LHD 是 $l{=}2$ heliotron 旋转椭圆、HSX 是准螺旋对称、CFQS 是低环径比准轴对称。这些需要完整的边界谐波。

\paragraph{做法。}
真机预设携带一个显式 \texttt{shape}：取\textbf{公开平衡}（DESC \texttt{desc/examples}：W7-X、HSX、ESTELL；经典 $l{=}2$ heliotron 作 LHD）的边界 Fourier 系数，\textbf{截断} $|m|,|n|\le2$、\textbf{归一化}（减 $R_{00}$、除原生小半径），按 DESC 双 Fourier 乘积基求值，再\textbf{重缩放}到本预设的 $R_0,a$：
\begin{equation}
R(\theta,\varphi)=R_0+a\!\sum_{m,n}\!R_{mn}A_m(\theta)B_n(\varphi),\quad
Z(\theta,\varphi)=a\!\sum_{m,n}\!Z_{mn}A_m(\theta)B_n(\varphi),
\end{equation}
\[
A_m(\theta)=\begin{cases}\cos m\theta & m\ge0\\ \sin|m|\theta & m<0\end{cases},\qquad
B_n(\varphi)=\begin{cases}\cos(n\nfp\varphi)& n\ge0\\ \sin(|n|\nfp\varphi)& n<0\end{cases}.
\]
这复现了 W7-X bean、LHD 旋转椭圆、HSX/CFQS 形态。\textbf{它是 cartoon}：$|m|,|n|\le2$ 截断、非完整平衡重建；进功率的 $\Vp,\Sw,\Sp$ 均由 \texttt{boundary\_metrics} 同一调用积分统一给出（$V_p$ 可用实测覆写，$\Sw,\Sp$ 不独立覆写）。

\codefrom{\texttt{stellarator.py} · \texttt{\_machine\_boundary\_outlines}（双 Fourier 乘积基求值）}
\begin{lstlisting}
# cos_m[m] = cos(m*theta) (m>=0) or sin(|m|*theta) (m<0)   # = A_m(theta)
Bn = math.cos(n*nfp*phi) if n >= 0 else math.sin(-n*nfp*phi)   # = B_n(phi)
Rh += c * _fade(m, rho) * cos_m[m] * Bn      # sum R_mn * A_m * B_n
# Z built the same way; final  R = R0 + a*Rh,  Z = a*Zh
\end{lstlisting}

%======================================================================
\section{体积与壁面：精确边界积分}\label{sec:integral}

无论概念堆还是真机，\textbf{功率账要的体积 $\Vp$、壁面 $\Sw$ 和等离子体表面积 $\Sp$ 都来自对边界的精确环向积分}（\texttt{boundary\_metrics} 同一调用中产出，$\Sp<\Sw$ 由构造保证），而且积分的就是\textbf{前端画出来的那一条边界}（前后端一致）。$\Vp$ 可用实测覆写替换；$\Sw$ 和 $\Sp$ 不独立覆写。

\paragraph{算法（\texttt{boundary\_metrics}）。}
边界 $\varphi$ 周期，所以只积\textbf{一个场周期} $[0,2\pi/\nfp)$ 再乘 $\nfp$（利用场周期对称性，对全周期积分验证一致到 $<10^{-13}$）。体积用 Green 定理把面积分化为线积分
\begin{equation}
\Vp=\oint\!\Big[\tfrac12 R^2\,dZ\Big]\,d\varphi\;\times\nfp,
\end{equation}
壁面（边界外扩间隙 $g$ 后）用叉积
\begin{equation}
\Sw=\oint\!\!\int\big|\partial_\theta \bm P\times\partial_\varphi \bm P\big|\,d\theta\,d\varphi\;\times\nfp,
\quad \bm P=(R\cos\varphi,\,R\sin\varphi,\,Z).
\end{equation}
对解析圆环面验证到 $0.004\%$。

\codefrom{\texttt{stellarator.py} · \texttt{boundary\_metrics}（体积，Green 定理）}
\begin{lstlisting}
phs = np.linspace(0.0, 2*math.pi/nfp, n_phi, endpoint=False)   # ONE field period
V = 0.0
for i in range(n_phi):
    R, Z = Rg[i], Zg[i]
    V += np.sum(0.5 * R * R * (np.roll(Z, -1) - Z)) * dph       # contour 1/2 R^2 dZ
V = abs(V) * nfp                                                # x field periods
\end{lstlisting}

\paragraph{概念堆怎么处理 $r^2$ 失效。}
概念堆积分的边界是\textbf{前端实际画的 bean}：满尺寸一阶体 + \textbf{限幅}二阶 wobble（$r^2$ 位移不超过 $r^1$ 的 $50\%$）。裸 $r^2$ 在 $r=a$ 会自交、体积暴涨（HELIAS 实测 $+331\%$）；限幅后是一条简单、尺寸正确的闭曲线，其积分比解析 Pappus 估计 $\pi a^2\Lax$ 低约 $2$--$6\%$（bean/曲率修正）。$\rho=1$ 边界与内面衰减律无关，故\textbf{功率账体积与显示嵌套画法解耦}。

\codefrom{\texttt{stellarator.py} · \texttt{\_nearaxis\_reconstruct}（限幅二阶重建）}
\begin{lstlisting}
# s1 = max |r1 displacement|, s2 = max |r2 displacement| at this phi
cap = min(a, _R2_DISPLAY_CAP * s1 / s2) if s2 > 0 else a    # _R2_DISPLAY_CAP=0.5
def surf(rho):
    wob = (rho ** _SHAPE_FADE) * cap * cap     # bounded r2; rho=1 -> cap^2
    R = R0p + rho * a * body_R + wob * wob_R   # full-size body + capped wobble
    Z = Z0p + rho * a * body_Z + wob * wob_Z
    return R, Z
\end{lstlisting}

\paragraph{面积锚。}
$\pi a^2$ 仍是一阶通量守恒截面积诊断（一阶椭圆面积恒为 $\pi a^2$，拉长比不改它）。求解器输出 $\Aflux=\Vp/\Lax$，并用 $a_\mathrm{vol}=\sqrt{\Vp/(2\pi^2R_0)}$ 进入回旋辐射、ISS04 和 Sudo。真机另带实测 $\Vp^{*}/\Sw^{*}$ 覆写。

%======================================================================
\section{几何怎么进功率账（哪个量、什么幂次）}\label{sec:account}

\begin{center}
\renewcommand{\arraystretch}{1.3}
\begin{tabular}{@{}lll@{}}
\toprule
功率量 & 公式（几何相关因子）& 几何依赖 \\
\midrule
聚变功率 & $\Pfus\propto n_{10}n_{20}\Phi\,\Vp$ & $\propto\Vp$ \\
轫致辐射 & $\Pbrem\propto n_e^2\sqrt{T_e}\,\Vp$ & $\propto\Vp$ \\
回旋辐射 & $\Pcycl\propto B_0^{2.5}a_\mathrm{vol}^{-0.5}\,\Vp$ & $\propto\Vp\,a_\mathrm{vol}^{-1/2}$ \\
储能/输运 & $\Eth\propto(n_iT_i+n_eT_e)\Vp$ & $\propto\Vp$ \\
第一壁负荷 & $\Pwall=(\Pfus+\Pheat)/\Sw$ & $\propto 1/\Sw$ \\
ISS04 约束 & $\tau\propto a_\mathrm{vol}^{2.28}R_0^{0.64}P_L^{-0.61}\bar n^{0.54}B_0^{0.84}\iotabar^{0.41}$ & $a_\mathrm{vol},R_0,\iotabar$ \\
Sudo 密度限 & $n_\mathrm{Sudo}=0.25\sqrt{P_LB_0/(a_\mathrm{vol}^2R_0)}\,10^{20}$ & $a_\mathrm{vol},R_0$ \\
\bottomrule
\end{tabular}
\end{center}

\begin{itemize}[leftmargin=1.4em]
\item \textbf{体积最强}：$\Pfus,\Pbrem,\Pcycl,\Eth$ 全线性正比 $\Vp$。真机用实测 $\Vp$，聚变/辐射锚在实测体积上。
\item \textbf{$\Sw$ 只进壁负荷}。$\Sw$ 和 $\Sp$ 均来自边界积分（同一调用，$\Sp<\Sw$ 由构造保证），不再独立覆写。
\item \textbf{$\iotabar$ 仅以 $0.41$ 次幂进 ISS04}。真机用实测 $\iotabar$；概念堆用近轴自洽解。
\item \textbf{形状经 $\Vp/\Sw$ 进功率（概念堆）}：功率公式无显式形状项，但 $\Vp,\Sw$ 是\textbf{所画 $\rho{=}1$ 边界的精确积分}（\S\ref{sec:integral}），所以改边界形状会经 $\Vp/\Sw$ 改功率。\textbf{实证}：NAE-QA 固定轴（$\Lax{=}114.1$）与 $a{=}1.8$，$\etabar{:}0.04\!\to\!0.07$（拉长比 $3.3\!\to\!4.5$）令 $\Vp{:}1095\!\to\!1145$、$\Pfus{:}1825\!\to\!1909$ MW（$\sim5\%$）。\textbf{真机}用实测 $\Vp/\Sw/\iotabar$ 覆写，所画形状不进功率（积分 $V_\mathrm{p,geom}$ 仅上报）。不进功率的是内面嵌套显示衰减（$\rho^{1.5}$）与 $B_{2c}$。
\end{itemize}

%======================================================================
\section{截面可视化：统一嵌套磁面 + 相位滑块}\label{sec:view}

\paragraph{统一嵌套画法。}
概念堆与真机\textbf{共用一套}嵌套构造：以磁轴（$m{=}0$ 谐波 / 近轴 $r\!\to\!0$ 极限）为芯，
\begin{equation}
S(\rho)=\text{axis}+\rho\,(\text{低阶体})+\rho^{1.5}\,(\text{高阶塑形}),\qquad \rho\in(0,1].
\end{equation}
低阶体（$m\!\le\!1$ 椭圆 / 一阶近轴）随尺寸线性缩放；高阶塑形（$m\!\ge\!2$ bean·triangle / 二阶 wobble）按 $\rho^{1.5}$ 衰减——\textbf{比原始近轴 $\rho^{|m|}{=}\rho^2$ 高一阶}。这一改动同时解决两个毛病：原 $\rho^2$ 衰减会让 W7-X $\varphi{=}0$ bean 尖端处 $\rho{=}0.85$ 面\textbf{穿出边界}（逐点 point-in-polygon 核验，$15/121$ 点在外）；$\rho^{1.5}$ 使磁面\textbf{严格嵌套不交叉}并向轴平滑圆化。$\rho$ 取均匀采样使径向间距均匀（W7-X $\varphi{=}0$ 间距不均匀度 $1.73\!\to\!1.11$，HELIAS $2.11\!\to\!1.15$），呈现 VMEC/NSE-global 那种平滑均匀套叠。$\rho=1$ 边界与衰减律无关，故\textbf{不影响第 \ref{sec:integral} 节的体积积分}。

\codefrom{\texttt{stellarator.py} · \texttt{\_fade}（真机谐波的 $\rho$ 衰减律）}
\begin{lstlisting}
def _fade(m, rho):
    am = abs(m)
    if am == 0: return 1.0      # m=0: axis position, fixed
    if am == 1: return rho      # m=1: ellipse body, linear in size
    return rho ** _SHAPE_FADE   # m>=2: bean/triangle shaping, _SHAPE_FADE=1.5
\end{lstlisting}

\paragraph{相位扫描。}
仿星器截面随 $\varphi$ 变，UI 提供\textbf{相位开关}（$\varphi=0/T4/T2$ 三个特征切面）与\textbf{连续相位滑块}（真机与概念堆\textbf{统一}支持——两条路径都输出一周期 120 帧，$\varphi{=}0$ 和 $\varphi{=}T$ 在滑块两端）。拖动时标签即时更新、重绘用 \texttt{requestAnimationFrame} 防抖；控件放在独立 \texttt{\#phControls} 容器中，不被 \texttt{drawShape()} 重建销毁，保证原生 drag 交互不被中断。

\paragraph{壁层与高级输入。}
壁层 $=$ 边界沿外法向外扩 $g$，保证 磁面 $\subset$ 边界 $\subset$ 壁。几何组顶部提供\textbf{三档输入模式}（默认近轴 / 磁轴 Fourier / 边界 Fourier），选定后只显示该路径真正需要用户编辑的参数；不适用的几何输入从 UI 消失，并在切换时从前端状态中清理，避免隐藏参数继续生效。前端 \texttt{clean()} 只透传数字、字符串、数组和白名单对象；当前唯一白名单对象是隐藏的 \texttt{shape}，用于 W7-X/LHD/HSX/CFQS 这类真机边界在编辑任意输入后仍随请求保留。清空真机边界的 $R/Z$ 输入才表示显式删除 \texttt{shape}、退回近轴路径。

%======================================================================
\section{预设库：9 个内置位形}\label{sec:presets}

\begin{center}
\renewcommand{\arraystretch}{1.25}
\begin{tabular}{@{}llccccccc@{}}
\toprule
预设 & 类型 & $\nfp$ & $R_0$[m] & $a$[m] & $B_0$[T] & $\iotabar$ & $\Vp$[m$^3$] & 备注 \\
\midrule
HELIAS     & 概念·近轴 & 5 & 18.0 & 1.80 & 5.0 & 0.82 & 1149 & HELIAS 反应堆族 \\
NAE-QA     & 概念·近轴 & 3 & 18.0 & 1.80 & 5.5 & 0.42 & 1145 & LSP r1 §5.1 缩放 \\
precise-QA & 概念·近轴 & 2 & 9.0  & 1.50 & 5.7 & 0.42 & 369  & Landreman--Paul 2022 \\
precise-QH & 概念·近轴 & 4 & 12.0 & 1.50 & 5.7 & 1.24 & 422  & Landreman--Paul 2022 \\
QH-nfp3    & 概念·近轴 & 3 & 12.0 & 1.50 & 5.7 & 1.25 & 311  & pyQSC 2022 QH nfp3 \\
\midrule
W7-X       & 真机·边界 & 5 & 5.5  & 0.55 & 2.5 & 0.88$^{*}$ & 30$^{*}$  & DESC W7-X，bean \\
LHD        & 真机·边界 & 10& 3.9  & 0.65 & 2.85& 0.40$^{*}$ & 30$^{*}$  & $l{=}2$ heliotron \\
HSX        & 真机·边界 & 4 & 1.2  & 0.15 & 1.0 & 1.05$^{*}$ & 0.4$^{*}$ & DESC HSX，QHS \\
CFQS       & 真机·边界 & 2 & 1.0  & 0.25 & 1.0 & 0.45$^{*}$ & 1.0$^{*}$ & ESTELL 代理，QA \\
\bottomrule
\end{tabular}
\end{center}
\footnotesize $^{*}$ 真机的 $\iotabar,\Vp$ 为\textbf{实测覆写}（进功率账）；$\Sw,\Sp$ 由边界积分统一给出。\normalsize

\paragraph{两类预设。}
\begin{description}[leftmargin=1.2em,style=nextline]
\item[概念堆（近轴定义）] HELIAS / NAE-QA / precise-QA / precise-QH / QH-nfp3。由近轴展开\emph{定义}（显式 $rc/zs$ + $\etabar$，无覆写），本程序\textbf{原生精确}算出几何与自洽 $\iotabar$。其中 precise-QA/QH/QH-nfp3 取自公开准对称近轴库（pyQSC / Landreman--Paul 2022），缩放规则 $rc=R_0\cdot rc_\mathrm{pub}$、$\etabar=|\etabar_\mathrm{pub}|/R_0$（与 NAE-QA 相同），$\iotabar$ \textbf{精确复现发表值}（QA 0.42、QH 1.24/1.25）。
\item[真机（边界谐波）] W7-X / LHD / HSX / CFQS。携带 $|m|,|n|\le2$ 归一化边界谐波，功率账用 $\iotabar$ 实测覆写和边界积分（$\Vp$ 可用实测覆写，$\Sw,\Sp$ 统一积分）。
\end{description}

%======================================================================
\section{POPCON 扫描与性能}\label{sec:popcon}

\paragraph{POPCON 是什么。}
在两个参数（默认 $T_{i0}\times n_{i0}$，$21\times21=441$ 点）的网格上，每点跑一次完整 0-D 求解，画出 $\Pfus,Q,\Pwall,\beta,H_\mathrm{ISS04}$ 等的等值线 + 最优运行窗——这就是\textbf{运行空间图（POPCON）}。

\paragraph{两个性能根因（均已无损修复）。}
\begin{enumerate}[leftmargin=1.5em]
\item \textbf{几何是扫描不变量、却每点重算}。POPCON 扫的 $T_{i0}/n_{i0}/B_0$ 都\textbf{不改几何}（几何只依赖 $R_0$、$a$、$\nfp$、$\delta_h$、$\etabar$、轴 $rc/zs$、$g$、边界 \texttt{shape}），但旧代码在 $441$ 个点 + 每次 $\tau_E/T_e$ 自洽递归都重算近轴解 + 边界积分。利用这个\textbf{场周期不变性}，把几何按几何参数 \emph{memoize}（\texttt{\_stellarator\_geometry}）：整张图只算\textbf{一次}几何；若扫描键\emph{是}几何量（$R_0/\etabar/\dots$）则键变、正确重算，永不 stale。

\codefrom{\texttt{stellarator.py} · \texttt{\_stellarator\_geometry}（按几何参数缓存）}
\begin{lstlisting}
key = (R0, a, int(round(N_fp)), delta_h, etabar, g,
       rc_key, zs_key, shape_key, B2c)    # GEOMETRY params only (no Ti0/ni0/B0)
cached = _GEOM_CACHE.get(key)
if cached is not None:
    return cached            # whole POPCON grid reuses ONE geometry
\end{lstlisting}
\item \textbf{numpy BLAS 多线程病态}。此 Windows numpy 2.4.6 的 OpenBLAS 对\textbf{中等矩阵}线程化异常：$121\times121$ 的 \texttt{linalg.solve} 多线程 $\sim0.6$ s、单线程 $\sim0.4$ ms（$\sim1300\times$），使近轴 $r^2$ 解（$\nphi{=}121$）单次 $\sim3.8$ s。numpy import 前 cap \texttt{OPENBLAS\_NUM\_THREADS=1}。无损：所有真实预设 $\Vp/\iotabar$ \textbf{逐位相同}，pyQSC 校验仍 $<10^{-9}$。

\codefrom{\texttt{app/server.py}（numpy 导入前 cap BLAS 线程）}
\begin{lstlisting}
for _v in ("OPENBLAS_NUM_THREADS", "OMP_NUM_THREADS", "MKL_NUM_THREADS",
           "NUMEXPR_NUM_THREADS"):
    os.environ.setdefault(_v, "1")   # MUST be before numpy loads OpenBLAS
import numpy as np
\end{lstlisting}
\end{enumerate}

\begin{center}
\renewcommand{\arraystretch}{1.25}
\begin{tabular}{@{}lcc@{}}
\toprule
 & 修复前 & 修复后 \\
\midrule
HELIAS 单次求解 & 4.07 s & 0.145 s \\
\textbf{HELIAS $21\times21$ POPCON} & \textbf{$\sim$29 分钟} & \textbf{0.17 s}（live 实测）\\
全测试套件 & $>4$ 分钟 & 34 s \\
\bottomrule
\end{tabular}
\end{center}
唯一代价：退化的\textbf{平面圆轴}测试（$\delta_h=0$，constant 曲率、零挠率，近轴方程病态）对 BLAS 求和顺序敏感 $\sim10^{-3}$，其 tokamak 对照 $\Vp$ 容差放宽 $5\!\times\!10^{-4}\to3\!\times\!10^{-3}$；真实位形不受影响。

%======================================================================
\section{一致性核查与诚实局限}\label{sec:consistency}

\paragraph{一致性（全部 9 预设）。}
\begin{itemize}[leftmargin=1.4em]
\item \textbf{核心一致}：前端 $a,R_0,\nfp$ $=$ 后端；一阶截面积诊断为 $\pi a^2$，求解器输出有效 $\Aflux=\Vp/\Lax$。
\item \textbf{概念堆}：前端所画近轴边界与后端积分为\textbf{同一条曲线}，体积/壁面\textbf{恒等一致}（$<10^{-6}$）；精确积分比 Pappus 估计低 $2$--$6\%$。
\item \textbf{真机}：功率账用体积/壁面积实测覆写（$>0$ 时替换 $\Vp/\Sw$）或边界积分；$\Sp$ 始终由边界积分给出；前端形状的 $\Vp^{\mathrm{geom}}/\Sw^{\mathrm{geom}}$ 作为诊断对比上报。
\item \textbf{新增 QS 预设}：原生精确复现发表 $\iotabar$（QA 0.42、QH 1.24/1.25）。
\end{itemize}

\paragraph{诚实局限。}
\begin{itemize}[leftmargin=1.4em]
\item 截面视图是 cartoon：真机用 $|m|,|n|\le2$ 截断 + 归一化重缩放，\textbf{非完整平衡重建}；LHD 用通用 $l{=}2$ heliotron、CFQS 用 ESTELL 代理（形态族正确、非逐机精确）。
\item 真机 $\Vp$ 覆写为公开估计（$\Sw,\Sp$ 由边界积分统一给出），需对照权威文献核实；预设中的 $\Sw^\mathrm{override}$ 仅作诊断对照（见代码 \texttt{\_notes}）。
\item 未新增\textbf{真机}预设：精确 DESC 双 Fourier 边界谐波需离线 DESC 管线，本环境不可复现，故只新增由近轴\emph{定义}的概念堆 QS 预设。
\item 嵌套面 $\rho\,$体$+\rho^{1.5}\,$塑形为显示近似（严格嵌套、均匀间距），但 $\rho{=}1$ 边界即功率账积分边界。
\end{itemize}

\paragraph{来源与代码。}
ISS04：Yamada \emph{NF} 45 (2005) 1684；Sudo \emph{NF} 30 (1990)；近轴：Garren \& Boozer \emph{PoF B} 3 (1991) 2805，Landreman--Sengupta--Plunk \emph{JPP} 85 (2019)；precise QS：Landreman \& Paul \emph{PRL} 128 (2022) 035001；边界谐波：DESC \texttt{desc/examples}（github.com/PlasmaControl/DESC）；近轴实现 pyQSC（github.com/landreman/pyQSC）。代码：\texttt{polyfusion/configs/stellarator.py}、\texttt{polyfusion/nearaxis.py}、\texttt{polyfusion/presets/stellarator.json}；几何$\to$功率细节见 docs/40，物理内核见 docs/39。

\end{document}
